% Zusammenfassung auf Deutsch, übersetzt aus sections/0-1_abstract_en.tex.
% Schweizer Schreibweise, kein Eszett. Aufbau wie im englischen Abstract: jeder
% Absatz beginnt mit der Frage, die der vorherige offen gelassen hat.
%
% Neu geschrieben am 2026-11-16. Die vorherige Fassung enthielt die Formulierung
% "bevor das Verfahren läuft", die CLAUDE.md 3.7 verbietet, nannte zwei
% Ergebnisse, die das Manuskript nicht mehr trägt, und beschrieb eine Arbeit,
% die nicht mehr Teil der Dissertation ist.
%
% Keine geprägten Begriffe, nach CLAUDE.md 3.5.

Ein Netzwerk, das zwei Zahlen modulo einer Primzahl addieren lernt und $2048$
Neuronen in seiner verborgenen Schicht hat, löst die Aufgabe fehlerfrei und
trägt die Berechnung auf etwa $332$ davon. Ändert man eine einzige Sache am
Training, nämlich welche Parameter sich im Aktualisierungsschritt einen
Normierungsfaktor teilen, steigt diese Zahl auf $2046$. Aufgabe, Daten,
Architektur und Breite bleiben gleich, und die Testgenauigkeit fällt von
$0.998$ auf $0.976$. Wie viel eines Netzwerks am Ende die Arbeit leistet, wird
also nicht dadurch entschieden, was das Netzwerk berechnen soll. Es wird
dadurch entschieden, wie das Training den Fehler unter den Neuronen aufteilt,
es lässt sich wählen, und die Wahl hat einen Preis.

Die Frage entstammt einer Übersicht über $87$ Agenten für die Computernutzung
und $33$ Datensätze. Das Feld nähert sich einem Rezept an, einem vortrainierten
Modell mit Bild- und Sprachfähigkeiten, das durch Prompting oder durch Klonen
menschlicher Demonstrationen angepasst wird, und sechs Einschränkungen treten
immer wieder auf. Vier betreffen, wie das Feld misst, was es baut, und unter
welchen Bedingungen es baut. Zwei betreffen das, was das Training
hervorgebracht hat: Diese Systeme verallgemeinern schlecht über die Situationen
hinaus, auf denen sie trainiert wurden, und benötigen grosse Mengen an
Demonstrationsdaten. Eine Übersicht kann diese beiden nicht erklären, weil sie
nie eine Aufgabe festhält und das Training variiert, und weil eine Erfolgsrate
zwei Netzwerke nicht trennen kann, die gleich abschneiden. Sie müssen im
Netzwerk selbst gemessen werden.

Dafür braucht es eine Aufgabe, deren Lösung Neuron für Neuron bekannt ist, und
die Arithmetik modulo einer Primzahl ist eine solche. Netzwerke, die darauf
trainiert werden, zeigen Grokking: Sie passen die Trainingsdaten fehlerfrei an,
bleiben über Tausende von Epochen auf Zufallsniveau bei ungesehenen Daten und
verallgemeinern dann plötzlich. Hält man Aufgabe, Daten, Architektur und Breite
fest und variiert nur die Regel, die den Fehler zuteilt, findet jede Regel, die
verallgemeinert, dieselbe Fourier-Konstruktion, und die Zahl der Neuronen, die
sie umsetzen, unterscheidet sich zwischen den Regeln um mehr als die Hälfte.

Eine Bedingung erklärt das. Ein Neuron bleibt in Gebrauch, wenn der Anteil am
Fehler, den es erhält, gross ist im Verhältnis zu jenem Teil seiner eigenen
Aktualisierung, der nicht aus dem Gradienten stammt. Das deckt drei
Mechanismen ab, für die es zuvor keine gemeinsame Erklärung gab: einen
Gewichtszerfallsterm, der zum Gradienten addiert wird, die numerische
Untergrenze im Nenner eines adaptiven Optimierers und den Zerfall eines
trainierbaren Verstärkungsfaktors. Nur ein Normierungsfaktor, den sich genau
ein Neuron teilt, kann dieses Verhältnis ändern, was die gebräuchlichen
Optimierer auf die eine Seite einer Grenze stellt und zwei Regeln auf die
andere. Gleicht man den Anteil am Fehler zwischen den Neuronen an, ohne die
Richtung eines einzigen Gradienten zu ändern, verschiebt sich die Berechnung
von $60$ auf $99$ Prozent der Schicht, zum Preis der oben genannten zwei
Punkte Genauigkeit. Gibt man ein Profil der Fehleranteile direkt vor und lässt
eine zufällige Permutation entscheiden, welches Neuron welche Grösse erhält,
verschiebt sich die Berechnung ebenso. Damit ist ausgeschlossen, dass der
Fehleranteil nur dem folgte, was die Neuronen ohnehin taten. Der Effekt wächst
mit der Breite der Schicht und verschwindet dort, wo das unveränderte Netzwerk
sie bereits ganz nutzt.

Die Wahl wird in beide Richtungen bezahlt. Eine über viele Neuronen verteilte
Berechnung verträgt das Entfernen von Neuronen deutlich besser. Eine auf wenige
Neuronen konzentrierte Berechnung lässt sich stärker beschneiden und passt sich
etwa doppelt so schnell an eine Aufgabe an, die neue Kombinationen der bereits
gelernten Merkmale verlangt, indem sie die ungenutzten Neuronen heranzieht.
Für eine Aufgabe, die nur ändert, welche Antwort jedes Merkmal stützt, gilt
dieser Vorteil nicht.

Die Lernregel ist nicht der einzige Teil des Trainings, der Struktur
hinterlässt. Hält man den Algorithmus zur Fehlerzuteilung fest und variiert das
Trainingsziel, so ist die Aktualisierung der Aufmerksamkeitsgewichte eines
Transformers eine Summe über die geordneten Token-Paare, die das Ziel zulässt.
Bidirektionales Training nimmt jedes Paar in beiden Reihenfolgen auf, und die
Gewichte bleiben nahezu symmetrisch, während autoregressives Training ein Token
einmal zum Ziel und vielfach zum Kontext macht, sodass einige wenige Spalten
dominieren. Die Herleitung bestätigt sich in gleich aufgesetzten Modellen mit
zwölf Schichten, und die Struktur zeigt sich in veröffentlichten Modellen über
Text, Bild und Audio hinweg. Eine Initialisierung, die aus der Herleitung
folgt, beschleunigt bidirektionales Training.

Eine Messung an einem trainierten Netzwerk taugt nur so viel wie das
Ausleseverfahren dahinter, auch die Messungen dieser Arbeit selbst. Das übliche
Ausleseverfahren für ein latentes Weltmodell wird auf kodierten Beobachtungen
angepasst und auf vorhergesagte angewandt, und wo der Trainingsverlust
Richtungen vergleicht und Beträge ignoriert, legt nichts den Betrag fest. Das
Verfahren berichtet dann vor allem dessen Abdriften: In synthetischen Daten
erreicht sein Fehler das $27$-fache dessen, der entsteht, wenn man stets den
mittleren Zustand vorhersagt, und bei einem veröffentlichten Modell entfernt
das Skalieren der Vorhersagen auf den Betrag der Kodierungen $97$ Prozent des
Fehlers, ohne etwas neu zu trainieren. Zur Korrektur gehören vier Prüfungen.
Wendet man sie auf das eigene Mass dieser Arbeit an, so zeigt sich, dass die
Richtung des zentralen Ergebnisses bestehen bleibt, während der Wert, den
dieses Mass angibt, der höchste von drei Massen ist.

Wenn sich dadurch entscheidet, wie viel eines Netzwerks eine Berechnung trägt,
und wenn das einen Preis hat, lässt es sich stattdessen in der Architektur
festlegen. Die letzte Studie stellt ein sensorisches Muster als eine Menge von
Neuronen dar, von denen jedes durch mehrere andere gehalten wird,
zusammengesetzt aus Fragmenten, die durch Hebbsche Plastizität gelernt wurden.
Rauschen, das nichts stützt, wird herausgefiltert, selbst wenn mehr als die
Hälfte der Eingabe verfälscht ist, teilweise verdeckte Figuren werden ergänzt,
und getrennt gelernte Fragmente setzen sich zu Mustern zusammen, die im
Training nie vorkamen. Das Modell lässt sich nicht über Strichzeichnungen
hinaus skalieren, und so schliesst die Arbeit damit, darzulegen, was eine
Darstellung dieser Art benötigen würde, um als Zustand des Umweltmodells eines
Agenten zu dienen.
